\begin{center}\begin{tikzpicture}[font=\small,>=stealth,line join=round]
\draw[gray!40,rounded corners=3pt] (0,0) rectangle (14.2,8.3);\node at (7.1,7.700000000000001) {例3：只有一个奇函数部分，整体变号};
\begin{scope}[shift={(4.1,3.9499999999999997)},x=1.0714285714285714cm,y=1.0454545454545452cm]\draw[->,gray] (-2.8,0)--(2.8,0) node[right]{$x$};\draw[->,gray] (0,-2.2)--(0,2.2) node[above]{$y$};\draw[blue,thick,solid,domain=-2.5:2.5,samples=180] plot(\x,{(\x)^3*cos(deg(\x))/(1+(\x)^2)});\end{scope}\node[anchor=west,align=left,text width=5.0cm] at (8.4,4.0) {$F=\dfrac{x^3\cos x}{x^2+1}$\\[12pt]分母恒正，定义域为 $\mathbb R$\\[10pt]奇 $\times$ 偶 $\div$ 偶\\[10pt]整体为奇函数};
\node[text width=12.6cm,align=center] at (7.1,.65) {曲线来自原函数；起伏不影响左右对应高度互为相反数。};
\end{tikzpicture}\end{center}
